\section{词向量的评估}

\subsection{内在评估与外在评估}

在 NLP 中，评估方法分为两大类：

\begin{importantbox}{两类评估方法}
\begin{itemize}
    \item \textbf{内在评估}（Intrinsic evaluation）：在特定的子任务上直接评估词向量质量（如类比测试、相似度评分）。优点：快速、帮助理解组件行为。缺点：与最终应用的相关性不确定。
    \item \textbf{外在评估}（Extrinsic evaluation）：将词向量嵌入到真实任务中（如命名实体识别、问答系统），评估任务性能的提升。优点：直接衡量实用价值。缺点：慢、难以隔离词向量本身的贡献。
\end{itemize}
\end{importantbox}

\subsection{内在评估：类比任务与相似度}

词向量的内在评估主要有两种方式：

\textbf{1. 向量类比任务}：形如 ``a : b :: c : ?'' 的类比题。例如 ``man : king :: woman : ?'' 应该得到 queen。通过在大量类比题上的正确率来评估词向量质量。

\begin{figure}[H]
\centering
\includegraphics[width=0.95\textwidth]{slides-images/slide-025.jpg}
\caption{GloVe 可视化：性别维度上的线性关系——man/woman、king/queen、brother/sister 等词对之间的方向近似平行\protect\footnotemark}
\end{figure}
\footnotetext{Slides 第26页。}

\textbf{2. 词相似度评分}：让人类标注者给词对打相似度分数（0-10 分），然后计算模型给出的余弦相似度与人类评分之间的\textbf{相关系数}。

\begin{knowledgebox}{人类相似度判断的数据集}
经典的人类相似度评估数据集包括 WordSim-353、SimLex-999 等。例如：
\begin{itemize}
    \item tiger/tiger: 10.0（完全相同）
    \item book/paper: 7.46（高度相关）
    \item plane/car: 5.77（同为交通工具）
    \item stock/phone: 1.62（关系较弱）
    \item stock/jaguar: 0.92（几乎无关）
\end{itemize}
好的词向量应该与人类判断高度相关。
\end{knowledgebox}

\subsection{外在评估：命名实体识别}

命名实体识别（Named Entity Recognition, NER）是一个经典的外在评估任务：识别文本中的人名、地名、组织名等实体。

将 GloVe 词向量加入到 NER 系统中，可以显著提升识别准确率——这说明词向量确实捕捉到了对下游任务有用的语义信息。

\subsection{本章小结}

评估词向量需要同时关注内在指标（类比正确率、与人类相似度判断的相关性）和外在指标（在下游任务上的性能提升）。内在评估帮助我们理解词向量的质量，外在评估告诉我们它们在实际应用中是否有用。GloVe 和 Word2Vec 在各类评估中表现相当，远优于简单的 SVD 方法。

%% ========== Part 5: 词义与多义性 ==========

